\section{The cyclic bar cycle and its relation}
\label{sh-section}

Two elementary homology inputs to the source's class construction
can be proved directly.  We give the actual cycle, compute the
homology group in which its diagonal images lie, and construct
a relation with a coefficient invertible modulo \(\ell^\nu\).
These give the cyclic homology calculation and the required
integral relation.\cite{alpo-uniform-source}

\subsection{The bar complex and the cyclic cycle}

For a group \(G\) and a commutative coefficient ring \(A\)
with trivial \(G\)-action, let \(C_j(G;A)\) be the free
\(A\)-module on the tuples
\([g_1|\cdots|g_j]\), with \(C_0(G;A)=A\).
For \(j\ge1\), put
\[
\begin{split}
 \partial[g_1|\cdots|g_j]
 ={}&[g_2|\cdots|g_j]\\
 &+\sum_{i=1}^{j-1}(-1)^i
       [g_1|\cdots|g_ig_{i+1}|\cdots|g_j]\\
 &+(-1)^j[g_1|\cdots|g_{j-1}].
\end{split}
\tag{SH1}
\]
In degree one the two empty tuples cancel.
The identity \(\partial^2=0\) follows by pairing the
two orders of deleting or multiplying any pair of
adjacent positions; the two terms have opposite signs.
The homology of this complex is \(H_j(G;A)\).
A homomorphism of groups acts on a tuple entry by entry
and commutes with \(\partial\).

Let \(q\ge2\), write \(C_q=\mathbb Z/q\mathbb Z\)
additively with generator \(t=\overline1\), and put
\(A=\mathbb Z/q\mathbb Z\).
For \(E\ge1\), define the following finite bar chain:
\[
 z_E=\sum_{j_1,\ldots,j_E\in C_q}
              [j_1|t|j_2|t|\cdots|j_E|t]
                   \in C_{2E}(C_q;A).
 \tag{SH2}
\]

\begin{lemma}\label{sh-cycle}
One has \(\partial z_E=0\).
\end{lemma}
\begin{proof}
The first face in (SH1) forgets \(j_1\).  Its sum over
\(j_1\) is therefore \(q\) times the same chain and is
zero in \(A\).
For \(1\le i<E\), pair the face multiplying
\(j_i\) and \(t\), with sign \(-1\), with the face
multiplying \(t\) and \(j_{i+1}\), with sign \(+1\).
After summing over all indices their unsigned values
are equal, since translation by \(t\) permutes \(C_q\)
in either variable.
Finally, the face multiplying \(j_E,t\), with sign
\(-1\), cancels the last face, with sign \(+1\),
by the same translation in \(j_E\).
These are all \(2E+1\) faces, proving the assertion.
\end{proof}

\subsection{The exact homology group}

\begin{proposition}\label{sh-cyclic-square}
For every \(m\ge0\),
\[
       H_m(C_q\times C_q;\mathbb Z/q\mathbb Z)
                    \simeq(\mathbb Z/q\mathbb Z)^{m+1}.
 \tag{SH3}
\]
\end{proposition}
\begin{proof}
Temporarily write \(t\) multiplicatively and put
\(B=A[C_q]\), \(N=1+t+\cdots+t^{q-1}\).
There is an augmented free \(B\)-resolution of the
trivial module \(A\):
\[
 \cdots\xrightarrow{\,t-1\,}B
       \xrightarrow{\,N\,}B
       \xrightarrow{\,t-1\,}B
       \xrightarrow{\,\varepsilon\,}A\longrightarrow0.
 \tag{SH4}
\]
To verify exactness, write \(u=\sum_{i=0}^{q-1}a_it^i\).
The equation \((t-1)u=0\) says that all \(a_i\) are
equal, so its kernel is \(AN=\operatorname{im}N\).
Also \(Nu=(\sum_i a_i)N\).
The kernel of \(N\) is therefore the same as the kernel
of augmentation.  The equations
\(b_{i-1}-b_i=a_i\), with cyclic indices, have a solution
precisely when \(\sum_i a_i=0\); successive partial
sums construct one.  Hence this kernel is
\((t-1)B\), as required.
These calculations hold over \(A\) without division.

The augmented resolution (SH4) is split exact as a
complex of \(A\)-modules.  Indeed \(AN\) is free of
rank one with a coordinate retraction, and the kernel
of augmentation is free of rank \(q-1\), with the
last coefficient determined by the others.
The short exact sequences just described consequently
split over \(A\).

Take the tensor product over \(A\) of two copies of
(SH4) and form the total complex.
It is a free resolution of \(A\) over
\[
           A[C_q]\otimes_A A[C_q]
                        =A[C_q\times C_q].
\]
For clarity, exactness follows by first taking homology
in one of the two nonnegative degrees: tensoring
preserves the explicitly split exact sequences above,
so only its degree-zero augmentation survives; the
other resolution then has homology \(A\) in degree
zero and zero otherwise.
There are exactly \(m+1\) summands
\(A[C_q\times C_q]\) in total degree \(m\).

After tensoring with the trivial module \(A\),
the maps \(t-1\) become zero and the maps \(N\)
become multiplication by \(q=0\) in \(A\).
The resulting complex has zero differentials and
degree-\(m\) term \(A^{m+1}\), proving (SH3).
This resolution computes the same homology as the
bar complex: lift the augmentation successively
between two free resolutions, and lift differences
of the resulting composites to construct chain
homotopies.  Thus their tensor complexes have
canonically isomorphic homology.
\end{proof}

\begin{lemma}\label{sh-primitive-relation}
Let \(A=\mathbb Z/\ell^\nu\mathbb Z\), where \(\ell\)
is prime and \(\nu\ge1\).
For any \(r+1\) vectors \(v_0,\ldots,v_r\in A^r\),
there are \(a_0,\ldots,a_r\in A\), at least one of
them a unit, with
\[
                         \sum_{i=0}^r a_iv_i=0.
\]
\end{lemma}
\begin{proof}
Lift the vectors to the columns of an integer
\(r\) by \(r+1\) matrix.
It has a nonzero rational kernel vector; clear
denominators and divide by the greatest common
divisor of its integer coordinates.
The resulting primitive integer vector is still
in the kernel.
At least one coordinate is not divisible by \(\ell\).
Reduce the relation modulo \(\ell^\nu\); that
coordinate becomes a unit.
The same proof works for \(r=0\), taking the
single coefficient to be \(1\).
\end{proof}

\begin{corollary}\label{sh-diagonal-relation}
Let \(q=\ell^\nu\), \(\nu\ge1\), and \(E\ge1\).
For \(0\le t\le2E+1\), define
\[
       d_t:C_q\longrightarrow C_q\times C_q,
                   \qquad x\longmapsto(x,tx).
\]
There are \(a_t\in\mathbb Z/q\mathbb Z\), an index
\(t_0\) with \(a_{t_0}\) a unit, and a bar chain
\(b\in C_{2E+1}(C_q\times C_q;\mathbb Z/q\mathbb Z)\)
such that
\[
                   \partial b
                      =\sum_{t=0}^{2E+1}a_t(d_t)_*z_E.
 \tag{SH5}
\]
\end{corollary}
\begin{proof}
All the displayed images are cycles by
Lemma~\ref{sh-cycle}.
Their \(2E+2\) homology classes belong to the module
\(A^{2E+1}\) computed in Proposition~\ref{sh-cyclic-square}.
Apply Lemma~\ref{sh-primitive-relation}.
The resulting zero homology class means precisely
that the indicated cycle is a boundary, providing \(b\).
\end{proof}

Under \(E=e+1\), the cycle (SH2) is exactly the
source's \(\mathrm{zc}(q,e)\), and the maps \(d_t\)
are its \(\mathrm{dt}(q,t)\).
Thus (SH5) has \(2e+4\) terms and proves the
actual relation required there.
The unused case \(\nu=0\) in the broad source
interface has the zero coefficient ring; all chains
are zero and the relation is immediate.

The prime-power restriction in
Lemma~\ref{sh-primitive-relation} cannot simply be
removed from its general module statement.
In \(\mathbb Z/6\mathbb Z\), a relation
\(3a+2b=0\) forces \(a\) even and \(b\) divisible
by \(3\), so neither coefficient can be a unit.
The source uses a prime-power modulus, as required.

These proofs establish the actual cyclic cycle and
relation, but do not yet supply the additive
finite-field cohomology functional or the decorated
arithmetic lattices that enter the later class
construction.


